\documentclass[]{article}

%opening
\title{Experience Based Discrimination: An Agent Based Model}
\author{}

\begin{document}

\title{Discrimination through Experience; A Contextual Bandit Problem} %%%%%%%%%%%%
\author{0783C}
\date{\today}
\maketitle

\let\thefootnote\relax
\footnotetext{Faculty of Economics, University of Cambridge} %%%%%%%%%%

\begin{abstract}
	This paper aims to construct an Agent Based Model of the labour market via a contextual bandit problem. The paper seeks to uncover mechanisms by which discrimination in the labour market occurs, and subsequently, policy that can remedy said mechanisms. While the model does not fully explain discrimination in the labour market, reflections upon it's results provide useful insight.
\end{abstract} %%%%%%%%%

\bigskip

\section*{Introduction}

Discrimination in the labour market is multi-faceted, and presents itself in many forms. Discrimination can be bare-faced in it's prejudice; it seeps in to the orifices of society, creating differences between groups. Yet these aren't the types of of discrimination usually referred to in economics. Within academic economics, the form of discrimination that matters it the one where individuals with the same productivity are paid differing wages.

\section*{Literature Review}

\subsection*{Discrimination}

LePage (2023) develops a model of Experience Based Dsicrimination. The paper, which largely motivates the dvelopment of the model presented in this paper, posits that firms are not only uncertain over individual productivities, but also the productivity of the group. TYpically economic letreature has compartmentalised discrimination into two main categories, \textit{taste-based}, and \textit{statistical} discrimination. The former, associated with Becker (1957), corresponds to discrimination purely because the hiring manager prefers (or dislikes) some group. The latter is predicated on the assumption that firms are not inherently discriminatory, but are discriminatory in practice due to failures in ability ot read signals. The literature is associated with the Arrow and Phelps lineage. Statistical Discrimination \cite[]{phelpsstat} arrives because of \textit{"scarcity of information about the existence and characteristics of workers and jobs"}. Thus, the 'Experience Based Discrimination' is closely related to the statistical discrimination literature. However, there is a large difference, as noted by LePage

\subsection*{Multi Armed Bandit Problems}

\bigskip

\subsection*{Reinforcement Learning}

\bigskip

\section*{Model}

Cowgill and Tucker (2019) note that the notion of bias, and thus outcomes are inherently linked to the objective function of the firm. To achieve equal productivity, that is equal payoff for the firm between groups, the firm must first decide their objective, hiring policy could seek to maximise efficiency, innovation, or some other internal objective idiosyncratic to some division of firm. Although the simplicity of the model makes it seem as though firms should just seek to maximise profit, the fact that it is a Multi Armed Bandit Model presents an additional metric to optimise: regret.

Regret is defined as the difference between the realised productivity of the arm chosen, and the productivity of the optimal arm (that which would be chosen by an oracle). 

$$R = P_c - P_{optimal}$$

In this model, it can be shown that minimising regret is equivalent to maximising profit. This motivates the use of BayesUCB to model soft (?) affirmative action policy - it aims to minimise regret by encouraging exploration of the arms, the data gathering allows firms to reduce regret (increase profit). While the firms objective function is essentially unchanged, the re-framing of the problem from one of profit maximisation to regret minimisation enlighten us that in cases where groups have the equal expected productivity, fairness and profit are equivalent objectives.
\section*{Simulation}

\bigskip

\section*{Discussion}

It should be of course acknowledged, that while each model purports some specific mechanism of discrimination, the true mechanism is likely to be some combination of each.

\begin{abstract}

\end{abstract}

\section{}

\end{document}
